Dos altaveus estan separats entre si una distància $2d$ i emeten un senyal harmònic i en fase. Produïm el so amb un generador de senyals que podem ajustar entre els 200 i els 400~Hz. La persona A està situada a $d_1 = 5,00\ \mathrm{m}$ de l'altaveu 1 i a $d_2 = 8,68\ \mathrm{m}$ de l'altaveu 2. Una segona persona B està situada a $d'_1 = 8,14\ \mathrm{m}$ de l'altaveu 1 i a $d'_2 = 4,00\ \mathrm{m}$ de l'altaveu 2.

\figura[0.5]{fig1.png}

\begin{apartat}{1,25}
Calculeu totes les freqüències que podem escollir en el generador perquè la persona A senti els senyals amb una intensitat màxima.\\
\hspace*{1.5em}Quina freqüència de les anteriors hem d'escollir perquè la persona A senti el senyal amb intensitat màxima i la persona B no el pugui sentir (és a dir, el senti amb intensitat mínima)?
\end{apartat}

\begin{apartat}{1,25}
Ara emetem música pels dos altaveus amb una potència de $2,36\times 10^{-3}\ \mathrm{W}$ cadascun. Calculeu la intensitat sonora generada per cada altaveu en la posició de la persona A. Determineu el nivell d'intensitat sonora que percep aquesta persona.
\end{apartat}

\nota{Considereu que les ones sonores es propaguen en les tres dimensions de l'espai i que la seva energia es distribueix en superfícies esfèriques.}
\dades{%
  $I_0 = 10^{-12}\ \mathrm{W\,m^{-2}}$.\\
  Velocitat del so en l'aire: $340\ \mathrm{m/s}$.\\
  Superfície esfèrica: $4\pi r^2$.}
